\clearpage\begingroup\SourceLeaf{43b}\HistoricalHead{38}{L}{PHYSIK. BRIEFWECHSEL.}
\begin{center}
[5.]\\[.7em]
\textsc{Gerling} an \textsc{Gauss}. Marburg, 6. October 1852.
\end{center}
\noindent .\ .\ .\ .\ .\ .\ .\ .\ .\ .\ .\ .\ .\ .\ .\ .\ .\ .\ .\ .\ .
\par\noindent\hspace*{\parindent}\{Mein desiderium hinsichtlich des \textsc{Foucault}schen Experiments besteht
aber noch fort. Ich denke mir nämlich unter dem Erdpol die Pendel-Ebene
im Raume ruhend und scheinbar auf dem Horizont bis in \(24^{\mathrm{h}}\) um \(360^{0}\) sich
fortdrehend, weil sich die Erde um die Verticale darunter wegdreht. —
Unter der Breite \(\varphi\) dreht sich scheinbar die Ebene (quasi-Ebene) des Pendelschwungs
um \(360^{0}\mathbin{.}\sin\varphi\). — Zwischen beiden Fällen besteht aber der Unterschied,
dass ich im ersten Falle mir evident \textls[120]{elementar} Rechenschaft geben
kann, warum es so seyn muss (weil die Schwingungs-Ebene, oder die gerade
Linie, welche den Bogen im tiefsten Punkt tangirt \textls[120]{constant im Raume}
bleibt) und dagegen im zweiten, da die Schwingungsebene theilweise
mitdreht, mir nicht beantworten kann: was bleibt hier constant, damit das
scheinbare Zurückbleiben der betreffenden Berührungslinie \(360^{0}\mathbin{.}\sin\varphi\) betragen
könne? — Ich weiss nicht, ob ich Ihnen erzählt oder schon einmal geschrieben
habe, dass die \textsc{Foucault}schen Experimente mich schon im vorigen Jahre veranlassten
zu fragen, ob nicht die in allen mir bekannten Fallversuchen vor-
\clearpage\SourceLeaf{44}\HistoricalHead{39}{R}{GERLING AN GAUSS.}
\noindent kommende Abweichung nach \textls[120]{Süden}, von welcher die Theorie dieser Versuche
keine Rechenschaft gibt, ihren Grund in der während der Zeit des Fallens
vorgehenden Drehung des auf dem Horizont beschriebenen Orientirungs-Kreuzes
haben könnte? Ich berechnete mir für die Fallzeit den Winkel dieser Drehung,
und multiplicirte mit dessen Tangente die östliche Abweichung, fand
aber die auf diese Weise berechnete südliche Abweichung viel kleiner (wenn
ich mich recht erinnere nur \(\tfrac{1}{5}\) bis \(\tfrac{1}{3}\) der beobachteten).\}\quad .\ .\ .\ .\ .\ .\ .\ .\ .\ .\ .\ .
\clearpage\endgroup
